% preambulo.tex — Formato del TP Final de Ingeniería Mecánica I (PDF).
% Manual de la cátedra: Times New Roman 12, alineación justificada,
% interlineado 1,15, A4. Lo incluye tp_mecanica/pdf.yaml.

\usepackage{graphicx}
\usepackage{pdfpages}          % carátula oficial (PDF de la carátula .docx completada)
\usepackage{pdflscape}         % láminas A3 reducidas, en página apaisada
\usepackage{float}
\floatplacement{figure}{H}
\usepackage{xurl}              % URLs largas de las notas: cortan en cualquier carácter
\usepackage{etoolbox}

% Interlineado 1,15 de Word. En Word, el renglón «sencillo» de Times New Roman
% 12 mide 13,8 pt (1,15 em); 1,15 × 13,8 = 15,87 pt. El renglón de LaTeX a
% 12 pt es de 14,5 pt: 15,87 / 14,5 = 1,0945.
\linespread{1.0945}

% Párrafos en bloque, separados 6 pt, como el informe anterior; justificado
% (alineación por defecto de LaTeX: no se usa \raggedright en el cuerpo).
\setlength{\parindent}{0pt}
\setlength{\parskip}{6pt plus 1pt}
\setlength{\emergencystretch}{2em}  % evita renglones desbordados sin romper la justificación

% Títulos: negrita, en la misma fuente, sin color
\usepackage{titlesec}
\titleformat{\section}{\fontsize{14}{17}\selectfont\bfseries}{}{0pt}{}
\titleformat{\subsection}{\normalsize\bfseries}{}{0pt}{}
\titleformat{\subsubsection}{\normalsize\bfseries\itshape}{}{0pt}{}
\titlespacing*{\section}{0pt}{18pt}{6pt}
\titlespacing*{\subsection}{0pt}{12pt}{4pt}
\titlespacing*{\subsubsection}{0pt}{10pt}{2pt}

% Índice con puntos guía también en las secciones
\usepackage{tocloft}
\renewcommand{\cftsecleader}{\cftdotfill{\cftdotsep}}
\renewcommand{\cftsecfont}{\normalfont}
\renewcommand{\cftsecpagefont}{\normalfont}
\setlength{\cftbeforesecskip}{3pt}
\renewcommand{\cfttoctitlefont}{\fontsize{14}{17}\selectfont\bfseries}
\AtBeginDocument{\tocloftpagestyle{tp}}   % sin esto, la página del índice sale con estilo plain (sin pie)

% Pie de página: el del informe anterior + número de página
\usepackage{fancyhdr}
\newcommand{\tppie}{Ing. Mecánica - Grupo 01 - Curso S1091 - Año 2026}  % lo redefine tp-mecanica.lua
\fancypagestyle{tp}{%
  \fancyhf{}%
  \fancyfoot[L]{\small\itshape\tppie}%
  \fancyfoot[R]{\small\thepage}%
  \renewcommand{\headrulewidth}{0pt}%
  \renewcommand{\footrulewidth}{0pt}}
\pagestyle{tp}

% Epígrafes de figuras y tablas: 11 pt; «Tabla», no «Cuadro»
\usepackage{caption}
\captionsetup{font=small,labelsep=colon,justification=centering}
\AtBeginDocument{\renewcommand{\tablename}{Tabla}}

% Tablas: 11 pt, filas algo más altas
\AtBeginEnvironment{longtable}{\small\renewcommand{\arraystretch}{1.25}}

% Notas al pie: 10 pt, justificadas, sangría francesa del número
\usepackage[hang,flushmargin]{footmisc}
\setlength{\footnotesep}{10pt}
\renewcommand{\footnotelayout}{\linespread{1.0}\selectfont}

% Viñetas: «•» de la fuente del texto, como en Word. babel-spanish pone las
% suyas al empezar el documento; esto va después.
\AddToHook{begindocument/end}{\renewcommand{\labelitemi}{\textbullet}\renewcommand{\labelitemii}{\textendash}}

% Bibliografía: sangría francesa de 1,27 cm (APA)
\newenvironment{tpbiblio}{\par\setlength{\parskip}{4pt}}{\par}
\newcommand{\tpitem}{\par\noindent\hangindent=1.27cm\hangafter=1 }
